\section{比較実験の設計}
\label{sec:experiments}

本節は評価の設計を述べる。
数値そのものは\cref{sec:results} に置き、
本節ではどの比較が何を切り分けるかを固定する。

\subsection{データ}
\label{sec:data}

\paragraph{合成側}
実会場の3DGS再構成から作ったデータセットを用いる。
\cref{tab:dataset_summary} に採用数を示す。
四つの会場 \SceneSet を対象とし、
trajectory group 単位で train / validation / test を分けている。
test 集合は trajectory-group-disjoint であり、
同じ会場の未学習視点を評価する。
これは未知会場への汎化ではない。

\begin{table}[t]
  \centering
  \caption{\textbf{合成データセットの採用数。}
  数値は生成 manifest（\texttt{dataset.json} と \texttt{splits.json}）から集計した。
  分割は trajectory group 単位で、
  同じグループのフレームは分割を跨がない。
  合計は公開記録にある \datasetEarlierReportTotal 枚と
  \datasetAcceptedBtwo の再生成差分だけ異なる。
  }
  \label{tab:dataset_summary}
  \small
  \begin{tabular}{@{}lrrrrr@{}}
    \toprule
    会場 & コート数 & 採用 & train & validation & test \\
    \midrule
    \texttt{B00} & 2 & \datasetAcceptedBzero & 1{,}718 & 220 & 238 \\
    \texttt{B01} & 3 & \datasetAcceptedBone & 1{,}559 & 295 & 198 \\
    \texttt{B02} & 1 & \datasetAcceptedBtwo & 1{,}648 & 194 & 252 \\
    \texttt{B03} & 1 & \datasetAcceptedBthree & 1{,}669 & 226 & 198 \\
    \midrule
    合計 & 7 & \datasetAcceptedTotal & \datasetTrainFrames &
      \datasetValidationFrames & \datasetTestFrames \\
    \bottomrule
  \end{tabular}
\end{table}


\paragraph{実写側}
\BaselineName 公開実装付属の実写データ
（\TCDImageCount 枚、十四点注釈）を用いる。
配布物は train と validation に分かれており、
本研究が確認した範囲では独立した test split が存在しない。
したがって実写側で報告できるのは held-out validation である。

\begin{table}[t]
  \centering
  \caption{\textbf{実写データの分割。}
  公開実装付属のデータには独立した test split が無い。
  したがって実写側で報告できるのは held-out validation であり、
  これを未知会場や非放送実写への汎化とは呼ばない。
  }
  \label{tab:data_partitions}
  \small
  \begin{tabular}{@{}lrrl@{}}
    \toprule
    集合 & 枚数 & 分割 & 本研究での扱い \\
    \midrule
    公開実写データ全体 & \TCDImageCount & --- & 十四点注釈 \\
    train & \TCDTrainCount & 75\% & 学習に使用 \\
    validation & \TCDValCount & 25\% & held-out validation \\
    test & --- & 無し & 未知会場・非放送の評価には別途データが必要 \\
    \bottomrule
  \end{tabular}
\end{table}


\subsection{比較相手}
\label{sec:baselines}

\paragraph{公開実装ベースライン}
\BaselineName\cite{tcdrepo} の固定 commit \BaselineCommit を用いる。
この実装は TrackNet 系のヒートマップ回帰\cite{huang2019tracknet}を
コート点検出に適用したもので、
補助中心点を含む十五チャネルを 640\(\times\)360 で出力する。
学習済み重みは公開配布物から取得し、
取得物の SHA-256 を付録~\ref{app:reproducibility} に記録する。
本研究は重みを再学習せず、
配布された検出器をそのまま比較相手として使う。

\paragraph{古典的手法}
Farin らの古典的コートモデル当てはめ\cite{farin2004courtmodels}の
派生実装\cite{gchlebusrepo}（commit \ClassicBaselineCommit、BSD-3-Clause）を
補助的な stress probe として走らせる。
主比較と同じ標本ではなく、
決定論的に選んだ実写 10 件と合成 \texttt{B00} 10 件を使う。
結果は\cref{sec:results_classical} に主比較とは別に示す。

\paragraph{比較手続き}
アーキテクチャも学習データも異なる検出器の生の点誤差を並べても、
幾何の良し悪しは決まらない。
そこで我々は、各検出器の生のキーポイントに対し、
\emph{同一の RANSAC ベースアライメント}を適用し、
その成功率と再投影残差を報告する。
これにより検出精度と登録精度を分ける。
公開実装は生のキーポイントからホモグラフィを解く段階も持つが、
確認した環境ではこの段階が現行の NumPy と SciPy で動作しないため、
本研究はこれを再現せず、生のキーポイントを共通手続きへ渡す。

共通手続きは正準規制テンプレートのキーポイントを画像へ写す
ホモグラフィを RANSAC で解く。
閾値は画像対角線の \resBenchRansacThreshold である。
有効な対応が足りない場合は失敗として数え、
成功率の分母から除かない。

\subsection{測定済みと未実施の境界}
\label{sec:measurement_boundary}

本稿の時点で測定済みなのは次の二つである。

\begin{enumerate}[leftmargin=1.6em]
  \item 実写 held-out validation と合成 trajectory-group test での
  \Method 対 \BaselineShort の交差モデル比較。
  domain ごとに決定論的 \resBenchDomainSamples 件、seed \resBenchSeed。
  \item 一つの checkpoint を選んだときに記録された
  合成 validation の値（\cref{sec:ckpt_validation}）。
\end{enumerate}

古典的コートモデル当てはめは、
主比較とは別の小標本での補助 stress probe までを測定済みとする。

\paragraph{まだ実施していない比較}
視点分布の効果を切り分けるため、次の条件を計画する。

\begin{table}[t]
  \centering
  \caption{\textbf{視点分布を切り分けるための比較計画。}
  本稿の時点ですべて未実施である。
  各条件は同じ再構成シーン、同じ教師 schema、同じモデル、
  同じ初期重み、同じ更新回数を持ち、
  動かす変数を一行に一つだけにする。
  }
  \label{tab:ablation_plan}
  \small
  \begin{tabularx}{\linewidth}{@{}p{0.05\linewidth}p{0.28\linewidth}X@{}}
    \toprule
    & 学習データ & 切り分けるもの \\
    \midrule
    A & 実写のみ & 実写学習の基準 \\
    B & 合成・視点狭域のみ（同枚数） & 制限された視点での基準 \\
    C & 合成・視点広域のみ（B と同枚数） & 視点範囲だけの効果 \\
    D & 同じ実写量 + 合成・狭域（同じ合成枚数） & 実写を使う条件の基準 \\
    E & 同じ実写量 + 合成・広域（D と同合成枚数） & 実写下での視点範囲の効果 \\
    \bottomrule
  \end{tabularx}
\end{table}


このうち A--E はチェックリストではなく、
本稿の主張に必要な順序である。
B 対 C が、同じ枚数で視点範囲だけを変えた比較、
D 対 E が、実写量と合成量を固定して視点範囲を変えた比較である。
この二つが揃って初めて、
\cref{sec:view_variable} の問いに答えられる。
したがって本稿の比較は、
視点拡張の効果そのものを分離しない。
アーキテクチャと学習データが同時に異なる条件での、
実測に基づく現状の性能差である。

採用された視点の分布そのものを描く集計は本稿には含めない。
視点の方位・高さ・距離の分布は\cref{sec:view_generation} の
生成契約から決まるが、
実際に学習へ渡った分布の集計図は未作成である。
本稿はこれを測定済みとして提示しない。

\subsection{指標}
\label{sec:metrics}

検出とアライメントの指標は、
次の定義で固定する。
数値は\cref{sec:results} に置き、定義はここに置く。

\paragraph{Complete（completeness）}
GT-visible なキーポイントのうち、
モデルが対応を返したものの割合である。
分母は GT-visible の点数であり、
位置誤差ではなく\emph{検出の網羅性}を測る。

\paragraph{PCK}
画像対角線に対する閾値
\(0.005, 0.01, 0.02, 0.05\) の割合で、
予測と GT の距離が閾値以内なら正解とする。
モデルが返さなかったキーポイントは不正解として数える。
したがって PCK は Complete と対応精度の両方に依存する。

\paragraph{pair error}
有効な対応だけを対象にした画素距離の平均・中央値・q90 である。
有効対応の件数を必ず併記する。
この統計は欠落を罰しないため、
PCK と Complete と一緒に読む。

\paragraph{共通ホモグラフィと成功率}
正準規制テンプレートから画像へのホモグラフィを
RANSAC で解き、その成功率を報告する。
分母は全標本であり、失敗も含める。
有効対応が不足する場合は失敗として扱い、
成功率から暗黙に除外しない。

\paragraph{線再投影誤差}
規制コートの線を各セグメント \resBenchLineSegmentSamples 点でサンプルし、
モデルのホモグラフィと GT ホモグラフィの対称再投影誤差を測る。
全サンプルの統計と、
GT の投影が画像内に落ちるサンプルだけの統計を分けて報告する。
前者は画面外への外挿を含むため、
撮影範囲の狭い標本では幾何の誤差より外挿距離を反映しやすい。
画面内率も併記し、
どれだけのコートがフレームに入っているかを示す。

\paragraph{コート多角形 IoU}
ダブルスコートの多角形を画像でクリップした後の IoU である。
クリップ後の多角形が縮退した標本には値が無い。
そのような標本を 0 として混ぜず、
定義できた件数と全件数を併記する。

\paragraph{姿勢指標}
カメラ中心の L2 距離（メートル）、
\(\SOthree\) の測地距離（度）、焦点距離の相対誤差、
対数焦点距離の絶対誤差、
および真値の正準点を予測姿勢で投影した再投影距離を使う。

再投影距離は\emph{姿勢のみ}の指標であり、
密なキーポイント予測を入力しない。
真のカメラ校正が無い実写画像に対して、
姿勢誤差を擬似的に作らない。
実写側では検出とアライメントの指標のみを報告し、
姿勢指標は合成 test に限る。
checkpoint 選択時の合成 validation が報告する姿勢指標は、
この定義に従う。

field registration の IoU と領域分割の mIoU は別の量として扱う。
前者はコートの平面的な重なり、
後者は画素単位の分類であり、
一つの数値へまとめない。

\paragraph{domain を跨ぐ比較の禁止}
実写と合成で GT-visible の定義が異なる。
実写は注釈が画像内にあること、
合成はカメラ前方・画面内・レンダラ可視の三条件である。
したがって本稿は、
Complete の絶対値を domain 間で比較しない。
domain 内のモデル比較は同じ分母で行う。

\subsection{失敗の扱い}
\label{sec:failures}

アライメントを棄却したシーンは結果から除かず、
negative evidence として報告する。
学習が完走しなかった run も、
可能な限り理由と停止時点を記録する。
失敗 run を結果を見た後に削除し、
成功した seed だけで平均を取る運用はしない。
補助 probe のクラッシュと縮退も同じ扱いである。
